\chapter{Proposition}
\label{chap:Propositions}

Partant de l'hypothèse qu'une représentation 3D pré-entraînée capture mieux les déterminants structurels de l'odeur qu'un GNN 2D appris à partir de zéro, nous proposons d'adapter le transformeur Uni-Mol à la classification multi-label des 138 descripteurs de GS-LF.
Cette proposition se décompose en trois sous-problèmes complémentaires~: la représentation des molécules et l'encodage par le transformeur, l'apprentissage du décodeur multi-label malgré le déséquilibre, et l'évaluation selon un protocole comparable à la littérature.
La section~\ref{sec:PropositionUniMol} présente successivement les idées qui la fondent, les outils nécessaires, sa formalisation et son analyse.

%--------------------------------------------------------------------------------

\section{Adaptation d'Uni-Mol à la classification multi-label des odeurs}
\label{sec:PropositionUniMol}

\subsection{Idées préliminaires}
\label{sec:IdeesPreliminaires}

\paragraph{Transférer plutôt qu'apprendre.}
Avec environ 5\,000 molécules, entraîner un transformeur de 47 millions de paramètres à partir de zéro conduirait à un surapprentissage massif.
Nous réutilisons donc Uni-Mol pré-entraîné comme \emph{encodeur}~: il transforme une molécule en un vecteur de dimension 512 (la sortie du token \texttt{[CLS]}), sur lequel un petit décodeur MLP prédit les 138 descripteurs.
Deux modes d'apprentissage sont envisagés~: encodeur entièrement gelé, où seul le décodeur apprend, puis dégel des dernières couches de l'encodeur, ajustées avec un taux d'apprentissage plus faible que celui du décodeur.

\paragraph{Utiliser la géométrie, et sa variabilité.}
Une molécule n'a pas une forme unique mais un ensemble de conformations.
Plutôt que d'en choisir une, nous en générons plusieurs par molécule et en tirons une au hasard à chaque passage pendant l'entraînement.
Ce tirage, complété par un léger bruit sur les coordonnées, constitue une augmentation de données naturelle~: le modèle apprend une représentation robuste aux variations de conformation, ce qui est précieux lorsque les données sont rares.

\paragraph{Compenser le déséquilibre.}
La plupart des descripteurs sont absents de la plupart des molécules~: une perte d'entropie croisée binaire standard est dominée par ces négatifs faciles, et le modèle apprend à prédire <<~absent~>> partout.
Nous comparons deux remèdes~: pondérer les positifs de chaque descripteur par l'inverse de leur fréquence, et utiliser une perte asymétrique \cite{Ridnik-2021}, qui atténue la contribution des négatifs faciles.

\paragraph{Évaluer comme la littérature.}
Les métriques n'ont de sens que relativement à un protocole.
Nous adoptons celui de \cite{Han-2026}, décrit à la section~\ref{sec:EAHan}, pour que nos résultats soient directement comparables à ceux de POM, OpenPOM et Uni-Mol2.

\subsection{Outils nécessaires}
\label{sec:OutilsNecessaires}

La proposition s'appuie sur les briques suivantes, détaillées au chapitre~\ref{chap:Experimentations}~:
\begin{itemize}
   \item \textbf{RDKit}\footnote{\url{https://www.rdkit.org}}, pour lire les codes SMILES, ajouter les hydrogènes et générer des conformations 3D~;
   \item \textbf{les poids pré-entraînés d'Uni-Mol} et la bibliothèque \texttt{unimol\_tools}, qui fournit l'architecture et le dictionnaire des atomes~;
   \item \textbf{PyTorch} et \textbf{PyTorch Lightning}, pour le décodeur, l'apprentissage et le calcul sur GPU~;
   \item \textbf{scikit-learn} et \textbf{scikit-multilearn}, pour les métriques et la stratification itérative multi-label \cite{Sechidis-2011, Szymanski-2017}.
\end{itemize}

\subsection{Formalisation}
\label{sec:Formalisation}

\subsubsection{Définition du problème}

Soit $\mathcal{D} = \{(m_i, \mathbf{y}_i)\}_{i=1}^{N}$ le jeu de données, où $m_i$ est une molécule décrite par son code SMILES et $\mathbf{y}_i \in \{0, 1\}^{L}$ le vecteur de ses descripteurs, avec $N = 4\,983$ et $L = 138$~: $y_{ij} = 1$ si le descripteur $j$ est attribué à la molécule $i$.
On cherche une fonction $f_\theta$ qui associe à une molécule un vecteur de probabilités $\hat{\mathbf{p}} = f_\theta(m) \in [0, 1]^{L}$, puis des seuils $\boldsymbol{\tau} \in [0, 1]^{L}$ tels que la prédiction $\hat{y}_j = \mathbb{1}[\hat{p}_j \geq \tau_j]$ soit la plus proche possible de $y_j$ pour chaque descripteur.

\subsubsection{Représentation 3D}

Pour chaque molécule $m$, on génère un ensemble de conformations $\mathcal{C}(m) = \{\mathbf{X}^{(1)}, \ldots, \mathbf{X}^{(K_m)}\}$, où $\mathbf{X}^{(k)} \in \mathbb{R}^{n \times 3}$ contient les coordonnées des $n$ atomes, hydrogènes compris.
Une conformation est encodée par~:
\begin{itemize}
   \item la suite des tokens $(t_0, t_1, \ldots, t_n, t_{n+1})$, où $t_0 = \texttt{[CLS]}$, $t_{n+1} = \texttt{[SEP]}$ et $t_1, \ldots, t_n$ sont les indices des atomes dans le dictionnaire d'Uni-Mol, de taille $V = 31$~;
   \item les coordonnées centrées $\mathbf{x}_k - \bar{\mathbf{x}}$, les tokens spéciaux étant placés à l'origine~;
   \item la matrice des distances $d_{kl} = \lVert \mathbf{x}_k - \mathbf{x}_l \rVert_2$~;
   \item le type de chaque paire, $e_{kl} = t_k \times V + t_l$, qui permet au modèle d'adapter l'encodage des distances aux éléments chimiques concernés.
\end{itemize}

\subsubsection{Encodeur}

Uni-Mol encode chaque distance par $K = 128$ noyaux gaussiens dont l'échelle et le décalage dépendent du type de paire~:
\begin{equation}
   \label{eq:Gaussiennes}
   \boldsymbol{\psi}_{kl} = \left( \mathcal{G}\left(a_{e_{kl}} \, d_{kl} + b_{e_{kl}} ;\ \mu_r, \sigma_r\right) \right)_{r = 1, \ldots, K},
   \qquad
   \mathcal{G}(x ; \mu, \sigma) = \frac{1}{\sigma \sqrt{2 \pi}} \, e^{-\frac{(x - \mu)^2}{2 \sigma^2}}
\end{equation}
Une projection linéaire de $\boldsymbol{\psi}_{kl}$ initialise la représentation de paires $\mathbf{q}^{0}_{kl}$, qui biaise l'attention de chaque couche $\ell$~:
\begin{equation}
   \label{eq:Attention}
   \mathrm{Attention}(\mathbf{Q}, \mathbf{K}, \mathbf{V}) = \mathrm{softmax}\left( \frac{\mathbf{Q} \mathbf{K}^{\top}}{\sqrt{d}} + \mathbf{q}^{\ell} \right) \mathbf{V},
   \qquad
   \mathbf{q}^{\ell + 1}_{kl} = \mathbf{q}^{\ell}_{kl} + \frac{\mathbf{Q}_k \mathbf{K}_l^{\top}}{\sqrt{d}}
\end{equation}
Après les 15 couches, la sortie du token \texttt{[CLS]} fournit la représentation de la molécule $\mathbf{h} = g_\phi(\mathbf{X}) \in \mathbb{R}^{512}$.
Les paramètres $\phi$ de l'encodeur sont gelés, à l'exception éventuelle des $u$ dernières couches.
Comme la représentation ne dépend que des distances, elle est invariante par toute isométrie de l'espace, rotations, translations \emph{et réflexions}.

\subsubsection{Décodeur multi-label}

Le décodeur est un perceptron multicouche à une couche cachée de dimension 256~:
\begin{equation}
   \label{eq:Decodeur}
   \mathbf{z} = \mathbf{W}_2 \, \mathrm{Dropout}\big(\mathrm{GELU}(\mathrm{LN}(\mathbf{W}_1 \mathbf{h} + \mathbf{b}_1))\big) + \mathbf{b}_2,
   \qquad
   \hat{\mathbf{p}} = \sigma(\mathbf{z})
\end{equation}
avec $\mathbf{W}_1 \in \mathbb{R}^{256 \times 512}$, $\mathbf{W}_2 \in \mathbb{R}^{138 \times 256}$, $\mathrm{LN}$ une normalisation de couche et $\sigma$ la fonction sigmoïde appliquée à chaque logit $z_j$ indépendamment.
Chaque ligne $\mathbf{w}_j$ de $\mathbf{W}_2$ est le vecteur propre au descripteur $j$~: le logit vaut $z_j = \mathbf{w}_j \cdot \mathbf{h}' + b_j$, où $\mathbf{h}'$ est la sortie de la couche cachée.
Les 138 descripteurs sont ainsi traités comme des tâches binaires indépendantes partageant la même représentation.

\subsubsection{Fonction de perte}

\paragraph{Entropie croisée pondérée.}
Pour un lot $\mathcal{B}$, la perte de référence est l'entropie croisée binaire dont les positifs du descripteur $j$ sont pondérés par $w_j = \min\left(\frac{N - N_j^{+}}{N_j^{+}},\ w_{\max}\right)$, où $N_j^{+}$ est le nombre de positifs du descripteur $j$ dans l'ensemble d'entraînement et $w_{\max} = 20$ un plafond qui évite des gradients excessifs pour les descripteurs les plus rares~:
\begin{equation}
   \label{eq:BCE}
   \mathcal{L}_{\mathrm{BCE}} = -\frac{1}{|\mathcal{B}|\,L} \sum_{i \in \mathcal{B}} \sum_{j=1}^{L} \Big[ w_j \, y_{ij} \log \hat{p}_{ij} + (1 - y_{ij}) \log (1 - \hat{p}_{ij}) \Big]
\end{equation}

\paragraph{Perte asymétrique.}
La perte asymétrique (ASL) \cite{Ridnik-2021} traite différemment positifs et négatifs.
Elle décale d'abord la probabilité des négatifs d'une marge $c$, $\hat{p}^{\,c}_{ij} = \max(\hat{p}_{ij} - c, 0)$, ce qui annule la perte des négatifs très faciles, puis les pondère par un facteur focal d'exposant $\gamma_-$~:
\begin{equation}
   \label{eq:ASL}
   \mathcal{L}_{\mathrm{ASL}} = -\frac{1}{|\mathcal{B}|} \sum_{i \in \mathcal{B}} \sum_{j=1}^{L} \Big[ y_{ij} \, (1 - \hat{p}_{ij})^{\gamma_+} \log \hat{p}_{ij} + (1 - y_{ij}) \, \big(\hat{p}^{\,c}_{ij}\big)^{\gamma_-} \log \big(1 - \hat{p}^{\,c}_{ij}\big) \Big]
\end{equation}
Nous utilisons $\gamma_+ = 0$, $\gamma_- = 2$ et $c = 0{,}05$~: les positifs, rares, conservent tout leur poids, tandis que les négatifs déjà bien classés contribuent peu.
Les facteurs focaux sont traités comme des constantes lors du calcul du gradient, comme dans l'implémentation de référence.

\subsubsection{Augmentation géométrique}

À chaque passage d'une molécule $m$ pendant l'entraînement, on tire une conformation $\mathbf{X}$ uniformément parmi les conformations relaxées de $\mathcal{C}(m)$, puis on perturbe ses coordonnées~:
\begin{equation}
   \label{eq:Jitter}
   \tilde{\mathbf{x}}_k = \mathbf{x}_k + \boldsymbol{\varepsilon}_k,
   \qquad
   \boldsymbol{\varepsilon}_k \sim \mathcal{N}(\mathbf{0}, s^2 \mathbf{I}_3),
   \qquad
   s \sim \mathcal{U}([s_{\min}, s_{\max}])
\end{equation}
avec $s_{\min} = 0{,}01$~\AA{} et $s_{\max} = 0{,}04$~\AA, des amplitudes faibles devant la longueur d'une liaison covalente (de l'ordre de 1~\AA).
En validation et en test, la conformation est fixe et aucun bruit n'est ajouté.
Une permutation aléatoire de l'ordre des atomes n'est pas utile~: l'attention sans encodage de position rend déjà Uni-Mol invariant à cet ordre.

\subsubsection{Protocole d'évaluation}

\paragraph{Découpage.}
GS-LF est découpé en un ensemble d'entraînement $\mathcal{D}_{\mathrm{train}}$ (80~\%) et un ensemble de test $\mathcal{D}_{\mathrm{test}}$ (20~\%) par stratification itérative d'ordre~2 \cite{Sechidis-2011, Szymanski-2017}, qui préserve la fréquence de chaque descripteur et de chaque paire de descripteurs.
$\mathcal{D}_{\mathrm{train}}$ est à son tour découpé en cinq plis $F_1, \ldots, F_5$ de la même manière.

\paragraph{Seuils hors pli.}
Pour chaque pli $k$, un modèle $f_k$ est entraîné sur $\mathcal{D}_{\mathrm{train}} \setminus F_k$, sa meilleure époque étant choisie selon l'AUROC macro sur $F_k$.
Chaque molécule d'entraînement reçoit ainsi une prédiction hors pli $\hat{\mathbf{p}}^{\mathrm{oof}}_i = f_{k(i)}(m_i)$, faite par le seul modèle qui ne l'a pas vue.
Le seuil de chaque descripteur maximise son F1 sur ces prédictions~:
\begin{equation}
   \label{eq:Seuils}
   \tau_j = \underset{\tau}{\arg\max} \ \mathrm{F1}\Big( \big(y_{ij}\big)_{i \in \mathcal{D}_{\mathrm{train}}},\ \big(\mathbb{1}[\hat{p}^{\mathrm{oof}}_{ij} \geq \tau]\big)_{i \in \mathcal{D}_{\mathrm{train}}} \Big)
\end{equation}

\paragraph{Ensemble et métriques.}
Sur le test, les probabilités des cinq modèles sont moyennées, $\hat{\mathbf{p}}_i = \frac{1}{5} \sum_{k=1}^{5} f_k(m_i)$, puis binarisées avec les seuils $\tau_j$.
Les métriques sont des moyennes macro, c'est-à-dire des moyennes non pondérées sur les descripteurs, ce qui donne le même poids aux descripteurs rares qu'aux fréquents~:
\begin{equation}
   \label{eq:F1Macro}
   \mathrm{F1}_{\mathrm{macro}} = \frac{1}{L} \sum_{j=1}^{L} \frac{2 \, P_j \, R_j}{P_j + R_j}
\end{equation}
où $P_j$ et $R_j$ sont la précision et le rappel du descripteur $j$.
L'AUROC et l'AUPRC macro, indépendantes des seuils, sont calculées de la même façon sur les descripteurs présents dans le test.
Enfin, l'incertitude est estimée en recalculant les métriques sur 1\,000 rééchantillonnages avec remise du test~; l'intervalle de confiance à 95~\% est donné par les percentiles 2,5 et 97,5.

\begin{figure}
   \centering
   \placeholder{schema: chaîne complète -- SMILES, RDKit (conformères), encodage Uni-Mol (tokens, distances, types de paires), augmentation, encodeur Uni-Mol, [CLS] 512, décodeur MLP, 138 logits, sigmoïde, seuils par descripteur}
   \caption{Vue d'ensemble de la chaîne proposée}
   \label{fig:Pipeline}
\end{figure}

\subsection{Analyse}
\label{sec:AnalyseProposition}

\paragraph{Ce que la proposition résout.}
Elle répond aux limites des GNN identifiées au chapitre~\ref{chap:EtatArt}~: l'attention globale capte les interactions entre atomes éloignés, la représentation de paires intègre la géométrie, et le pré-entraînement compense en partie la rareté des données.
La perte asymétrique et les seuils par descripteur traitent le déséquilibre à deux niveaux distincts, l'apprentissage et la décision.
Le protocole d'évaluation rend la comparaison avec la littérature possible.

\paragraph{Ce qu'elle ne résout pas.}
L'invariance par réflexion (équation~\ref{eq:Attention} et suivantes) empêche toute distinction entre énantiomères, et GS-LF, dépourvu de stéréochimie, ne permettrait de toute façon pas de l'apprendre.
Les descripteurs restent traités comme indépendants, alors que leurs co-occurrences sont fortes et que les confusions du modèle en dépendent \cite{Han-2026}.
Enfin, la variabilité des annotations humaines borne la performance atteignable, quel que soit le modèle.

\paragraph{Coût de calcul.}
L'attention est quadratique en nombre d'atomes~: avec les hydrogènes, les molécules de GS-LF comptent jusqu'à 184 atomes.
Le dégel de couches impose de rétropropager à travers l'encodeur, ce qui contraint la taille des lots~; nous utilisons des lots de 8 avec accumulation de gradient sur 8 lots.
Le protocole complet multiplie le coût d'un entraînement par cinq, contre cinquante pour celui de Han et al.

%--------------------------------------------------------------------------------

\section{Conclusion}
\label{sec:ConclusionProposition}

Notre proposition combine un encodeur 3D pré-entraîné, une augmentation par conformations, un décodeur MLP entraîné avec une perte asymétrique et un dégel partiel de l'encodeur, ainsi qu'un protocole d'évaluation aligné sur la littérature.
Plusieurs de ses composantes sont des choix à valider empiriquement~: le mode d'apprentissage, la fonction de perte, l'intérêt de l'augmentation, la profondeur du décodeur ou la stratégie de découpage.
Le chapitre suivant décrit leur mise en œuvre et les expériences menées pour les départager.
